% Source coverage: physical PDF pp. 168--171 (printed pp. 145--148), beginning at the optional-reading Section 18.5 heading on p. 168 and ending with the universality-class paragraph immediately before the Section 18.6 heading on p. 171.
% Numbered equations: (18.5.1)--(18.5.13), with no unnumbered displays. Figures/tables: none. Footnotes: 4. Inline chapter references: [7], [8a], [8].
% Source review: rendered physical pp. 168--171 checked, including both sides of the Section 18.5/18.6 boundary. Corrected errata: T_0 is T_c in Eq. (18.5.3) and its first body footnote; the remainder in Eq. (18.5.8) begins at O(g_4^2); and the closing expansion is in powers of \epsilon.

\subsection[Critical Phenomena]{Critical Phenomena\footnotemark}
\label{sec:18.5}
\footnotetext{This section lies somewhat out of the book's main line of development, and may be omitted in a first reading.}

For some purposes we may be interested in the limit of very low rather than high energies or wave numbers. The arguments of Section 18.2 can be repeated to study this limit, except that here we must examine the case $\mu\to0$ rather than $\mu\to\infty$. This limit is of course simplest if there are no masses in the theory, as, for instance, in quantum electrodynamics with a symmetry under the chiral transformation $\psi\to\gamma_5\psi$ which forbids an electron mass. In this particular case the only renormalizable coupling $eA^\mu\bar\psi\gamma_\mu\psi$ as well as all non-renormalizable couplings have $\beta^\ell>0$ for sufficiently small couplings, so all trajectories in at least a finite region around the origin are attracted into the point $g^\ell=0$ as $\mu\to0$.

The same considerations may be applied even to theories with very small but non-zero masses if we include these masses among the coupling parameters of the theory, as described in the previous section. The coefficient $\Delta$ in Eq.~\eqref{eq:18.4.12} is positive for a mass parameter, so in this case the trajectories can never reach the point $g=0$, but they may come close if the masses are small.

Of course, even if we can regard some degree of freedom like the electron field as having zero or very small mass, in the real world there are many other degrees of freedom whose masses are not small. The renormalization group should properly be applied not to the true theory that encompasses all these heavy degrees of freedom, but to an `effective' field theory, in which only massless or nearly massless degrees of freedom appear explicitly, with interactions that include the effects of internal heavy particle lines. (We shall have more to say about effective field theories in Chapter 19.)

The low wave number limit is of particular interest in the study of critical phenomena, such as long-range correlations at or near a second-order phase transition (a smooth phase transition, with no latent heat) in condensed matter. Because we are interested in the limit $\mu\to0$, the important eigenvectors of the matrix \eqref{eq:18.4.4} are those with eigenvalues $\lambda<0$, which are called \emph{relevant}. The eigenvectors with $\lambda=0$ and $\lambda>0$ are called \emph{marginal} and \emph{irrelevant}, respectively.

Suppose that there is a non-trivial fixed point $g_*$ with just one negative eigenvalue $\lambda_0$, perhaps corresponding approximately to a mass operator. The set of trajectories of $g^\ell(\mu)$ that are attracted into this fixed point for $\mu\to0$ therefore forms a critical surface of codimension one; that is, a surface defined by a single condition on the couplings, the condition that for $g\to g_*$, the tangents $g^\ell-g_*^\ell$ have no components in the direction of the eigenvector with negative eigenvalue. There is a phase transition as the physical value of the couplings at any fixed characteristic scale approaches this surface. Because the critical surface has codimension one, the phase transition can be reached by adjusting any one parameter on which the couplings depend, such as the pressure or temperature. The fact that a wide variety of substances do exhibit phase transitions of this sort shows that it is common to encounter fixed points for which the matrix \eqref{eq:18.4.4} has a single negative eigenvalue, as already mentioned in the previous section.

To be specific, as the temperature $T$ approaches its critical value $T_c$ we expect the coefficient $c_0$ of the growing term in Eq.~\eqref{eq:18.4.5} to become proportional to $T-T_c$, because there is no reason why it should be singular or why it should vanish faster than this. Hence for $\mu\to0$ and then $T\to T_c$, the couplings go as
\begin{equation}
g^\ell(\mu)\to(T-T_c)V_0^\ell\mu^{\lambda_0},
\tag{18.5.1}\label{eq:18.5.1}
\end{equation}
where $\lambda_0$ is the only negative eigenvalue at $g_*$, and $V_0^\ell$ is the corresponding eigenvector.\footnote{Other contributions to the couplings will go as $(T-T_c)^0\mu^{\lambda_1}$ with $\lambda_1>0$. Thus Eq.~\eqref{eq:18.5.1} is valid here provided $T-T_c$ does not go to zero as fast as $\mu^{\lambda_1-\lambda_0}$.} Applying our renormalization group arguments to wave numbers instead of energies, the $N$-point function (the $N$th partial derivative of the effective action with respect to a field $\phi$ of dimensionality $[\text{wave number}]^{D_\phi}$) at a small characteristic wave number scale $\kappa$ has the form\footnote{The function $F_N$ also depends on dimensionless angles and ratios of wave numbers. Note that $\Gamma_N$ has `naive' dimensionality $d-ND_\phi$ because $\delta^d(\sum\kappa)\Gamma$ must be dimensionless.}
\begin{equation}
\Gamma_N(\kappa)\to\kappa^{d-N[D_\phi+\gamma_\phi(g_*)]}
F_N\bigl((T-T_c)\kappa^{\lambda_0}\bigr),
\tag{18.5.2}\label{eq:18.5.2}
\end{equation}
where $\gamma_\phi(g)$ is the anomalous dimension associated with the field $\phi$, and $d$ is the spacetime dimensionality, or in classical statistical mechanics the spatial dimensionality. It is convenient to rewrite this in the equivalent form
\begin{equation}
\Gamma_N(\kappa)\to(T-T_c)^{-[d-N(D_\phi+\gamma_\phi(g_*))]/\lambda_0}
G_N\bigl(\kappa(T-T_c)^{1/\lambda_0}\bigr).
\tag{18.5.3}\label{eq:18.5.3}
\end{equation}
This shows for one thing that the correlation length $\xi$ (the characteristic length that determines the scale over which the Fourier transform of $\Gamma_N$ varies) increases as $T$ approaches $T_c$ like
\begin{equation}
\xi\propto(T-T_c)^{-\nu},
\tag{18.5.4}\label{eq:18.5.4}
\end{equation}
where $\nu$ is a conventionally defined positive `critical exponent', given by Eq.~\eqref{eq:18.5.3} as
\begin{equation}
\nu=-1/\lambda_0.
\tag{18.5.5}\label{eq:18.5.5}
\end{equation}

Also, the zero-field effective action $\Gamma_0$ (or in statistical physics, the free energy) must be $\kappa$-independent because it corresponds to graphs with no external lines. It follows that Eq.~\eqref{eq:18.5.3} here becomes for $T\to T_c$:
\begin{equation}
\Gamma_0-F_0\propto(T-T_c)^{\nu d},
\tag{18.5.6}\label{eq:18.5.6}
\end{equation}
where the constant $F_0$ is the effective action or free energy due to the heavy degrees of freedom which have been integrated out. Thus the exponent $\nu$ also governs the behavior of the part of the free energy which is not analytic in temperature for $T$ near $T_c$.

In 1972 Wilson and Fisher~\hyperref[ch18-ref-7]{[7]} used an expansion in powers of $d-4$ both to show that the theory of a scalar field actually fits the above description, and also to carry out an approximate calculation of critical exponents such as $\nu$. Consider a theory with a single `light' degree of freedom, a scalar field $\phi$, such as the magnetization in a ferromagnet, with a symmetry under $\phi\to-\phi$ that rules out interactions odd in $\phi$. In addition to the `mass' term $-g_2\phi^2/2$, the Lagrangian density of the effective field theory will contain interactions $-g_4\phi^4/4!$, $-g_6\phi^6/6!$, etc. The dimensionality of the field $\phi$ in powers of wave number is $(d-2)/2$ (so that $\int d^dx\,(\nabla\phi)^2$ should be dimensionless) so the dimensionalities of the couplings $g_2$, $g_4$, $g_6$, etc. in $d$ dimensions are $+2$, $4-d$, $6-2d$, etc. For the fixed point at zero coupling in three dimensions, there are two relevant couplings, $g_2$ and $g_4$, but this conclusion is changed by interactions at non-trivial fixed points. Let's examine the surface in coupling constant space in which only $g_2$ and $g_4$ are non-zero, and take $g_4$ to be small.\footnote{These are the only renormalizable couplings for $3\le d\le4$, so for such $d$ this is an invariant surface. Note that we do not include the coefficient of $(\nabla\phi)^2$ among the couplings here, because this is a redundant coupling, in the sense described in Section 7.7.} Eq.~\eqref{eq:18.2.12} gives $\beta(g_4)=3g_4^2/16\pi^2+O(g_4^3)$ for $d=4$, and Eq.~\eqref{eq:18.4.12} tells us that for $d=4-\epsilon$ dimensions we must add to this a term $-\epsilon g_4$, so
\begin{equation}
\mu\frac{d}{d\mu}g_4(\mu)=-\epsilon g_4(\mu)
+\frac{3g_4^2(\mu)}{16\pi^2}+O\bigl(g_4^3(\mu)\bigr).
\tag{18.5.7}\label{eq:18.5.7}
\end{equation}
Also, Eq.~\eqref{eq:18.4.17} gives
\begin{equation}
\mu\frac{d}{d\mu}g_2(\mu)
=\left[-2+\frac{g_4(\mu)}{16\pi^2}+O\bigl(g_4^2(\mu)\bigr)\right]g_2(\mu).
\tag{18.5.8}\label{eq:18.5.8}
\end{equation}
Therefore for small $\epsilon$ there is a non-trivial fixed point at
\begin{equation}
g_{4*}=\frac{16\pi^2\epsilon}{3},\qquad g_{2*}=0.
\tag{18.5.9}\label{eq:18.5.9}
\end{equation}

The matrix \eqref{eq:18.4.4} at this fixed point is diagonal, with eigenvalues
\begin{equation}
\lambda_4=M^4{}_4=-\epsilon+\frac{3g_{4*}}{8\pi^2}
+O(g_{4*}^2)=+\epsilon+O(\epsilon^2),
\tag{18.5.10}\label{eq:18.5.10}
\end{equation}
\begin{equation}
\lambda_2=M^2{}_2=-2+\frac{g_{4*}}{16\pi^2}
+O(g_{4*}^2)=-2+\frac{\epsilon}{3}+O(\epsilon^2).
\tag{18.5.11}\label{eq:18.5.11}
\end{equation}
From Eq.~\eqref{eq:18.5.10} we see that the coupling $g_4$ is actually irrelevant, so that there is just one relevant coupling here, signalling the presence of a second-order phase transition. From Eq.~\eqref{eq:18.5.11} we see that the anomalous exponent \eqref{eq:18.5.5} is
\begin{equation}
\nu=-\frac{1}{\lambda_2}=\frac{1}{2}+\frac{\epsilon}{12}+O(\epsilon^2).
\tag{18.5.12}\label{eq:18.5.12}
\end{equation}
For the physical value $\epsilon=1$ the first two terms give $\nu\simeq0.58$. Three-loop calculations~\hyperref[ch18-ref-8a]{[8a]} give this critical exponent to order $\epsilon^3$ as
\begin{equation}
\nu=\frac{1}{2}+\frac{\epsilon}{12}+\frac{7\epsilon^2}{162}
-0.01904\epsilon^3,
\tag{18.5.13}\label{eq:18.5.13}
\end{equation}
which for $\epsilon=1$ gives $\nu=0.61$.

In the calculation presented here nothing was assumed about the system under study except that there is a second-order phase transition, near which the only long-wavelength degree of freedom is a single scalar field. There are a number of different physical systems that fit this description, such as the spontaneous appearance of magnetization (represented here by $\phi$) in ferromagnetic and antiferromagnetic materials, and also second-order phase transitions between liquids and gases and in binary fluids. All of these systems are therefore expected to have the same value of $\nu$. This is confirmed by experiment, which gives a value~\hyperref[ch18-ref-8]{[8]} $\nu=0.63\pm0.04$, in good agreement with the three-loop result \eqref{eq:18.5.13}, and in fair agreement even with the one-loop result \eqref{eq:18.5.12}. It is fortunate though still somewhat mysterious that an expansion in powers of $\epsilon$ should work so well.

More generally, all the systems that are described by the same set of long-wavelength degrees of freedom near their second-order phase transitions are said to belong to the same universality class. All the critical exponents are the same for all systems in a given universality class.
